\chapter{Conception}
\begin{spacing}{1.2}
\minitoc
\thispagestyle{MyStyle}
\end{spacing}
\newpage

{\Large\bfseries Introduction}\\
Après l’analyse des besoins, ce chapitre présente la phase de conception. Nous utilisons UML
pour modéliser l’architecture du système. Notre démarche inclut : diagrammes de cas d’utilisation,
de classes, de séquence et d’activité. Ces modèles guideront la phase d’implémentation.
\section{Choix de la méthodologie d'analyse}
Pour structurer notre projet \reportAppName{}, nous utilisons UML, un langage de modélisation
orienté objet. Il offre une panoplie de diagrammes statiques et dynamiques pour une représentation
précise du système.
\subsection{Présentation UML}
UML (Unified Modeling Language) est un langage de modélisation visuel standardisé, essentiel
en génie logiciel pour visualiser, construire et documenter les systèmes orientés objet. Il offre
un ensemble de diagrammes (cas d’utilisation, classes, séquences, activités) permettant de
modéliser la structure statique et le comportement dynamique d’un logiciel, facilitant la communication
entre les parties prenantes du projet.
\subsection{Description détaillée du diagramme de cas d'utilisation : }
Les diagrammes suivants détaillent les fonctionnalités du système du point de vue des utilisateurs. Ils
montrent comment les acteurs interagissent avec l'application \reportAppName{}.
\subsubsection{Diagramme de cas d'utilisation « Consulter le site » :}
Ce cas permet au visiteur de consulter le site ou l’application \reportAppName sans authentification.
Il offre une présentation des services et produits de \reportEntreprise{}.\\

\begin{figure}[H]%
 \center%
 \setlength{\fboxsep}{0pt}%
 \setlength{\fboxrule}{0.0pt}%
 \fbox{
 \includegraphics[width=\linewidth]{consulter_le_site.png}%
 }
 \caption{Diagramme de cas d'utilisation « Consulter le site »}%
\end{figure}

\textbf{Tableau descriptif :}
\begin{table}[h!]
\renewcommand{\arraystretch}{1.8}
\setlength{\tabcolsep}{16pt}
\centering
\begin{tabular}{|l|p{11cm}|}

\hline
\textbf{Titre} & \textbf{Consulter le site} \\
\hline
Objectif & Permettre au Visiteur de consulter les informations publiques de \reportAppName{}. \\
\hline
Acteur(s) & Visiteur \\
\hline
Pré-condition & Aucune, accès public sans authentification. \\
\hline
Post-condition & Le Visiteur a consulté les informations et peut s'authentifier. \\
\hline
Scénario nominal & \begin{itemize}
 \item Le Visiteur accède à l'URL de l'application \reportAppName{}
 \item Le système affiche la page d'accueil avec les informations publiques
 \item Le Visiteur navigue entre les différentes sections publiques
 \item Le Visiteur peut accéder à la page de connexion pour s'authentifier
\end{itemize}
\\
\hline
Scénario alternatif & Aucun — accès public sans restriction. \\
\hline
\end{tabular}
\textbf{\caption{Tableau descriptif du cas d'utilisation « Consulter le site »}}
\label{table:tableau descriptif du cas d'utilisation « Consulter le site »}
\end{table}
\vfill
\subsubsection{Diagramme de cas d'utilisation « sauthentifier » :}
Ce diagramme détaille l'authentification des utilisateurs. Le cas « S'authentifier » est étendu par « Rediriger selon rôle » vers le tableau de bord approprié.\par

\begin{figure}[H]%
 \center%
 \setlength{\fboxsep}{0pt}%
 \setlength{\fboxrule}{0.0pt}%
 \fbox{
 \includegraphics[width=\linewidth]{sauthentifier.png}
 }
 \caption{Diagramme de cas d'utilisation « sauthentifier »}%
\end{figure}

\textbf{Tableau descriptif :}
\begin{table}[!htbp]
\renewcommand{\arraystretch}{1.8}
\setlength{\tabcolsep}{16pt}
\centering
\begin{tabular}{|l|p{11cm}|}

\hline
\textbf{Titre} & \textbf{S'authentifier} \\
\hline
Objectif & Permettre à tout utilisateur de s'authentifier et d'accéder à son tableau de bord. \\
\hline
Acteur(s) & Client, Responsable Ventes, Magasinier, Responsable Achats, Responsable Production, PDG \\
\hline
Pré-condition & L'utilisateur possède un compte actif dans le système. La page de connexion est accessible. \\
\hline
Post-condition & L'utilisateur est authentifié et redirigé vers le tableau de bord correspondant à son rôle. \\
\hline
Scénario nominal & \begin{itemize}
 \item L'utilisateur accède à la page de connexion
 \item L'utilisateur saisit ses identifiants (Login et Mot de passe)
 \item Le système valide les informations
 \item Le système redirige l'utilisateur vers le tableau de bord approprié selon son rôle (Rediriger selon rôle)
\end{itemize}
\\
\hline
Scénario alternatif & Identifiants invalides : \begin{itemize}
 \item Le système affiche un message d'erreur « Login ou Mot de passe incorrect »
 \item L'utilisateur peut réessayer ou utiliser « Mot de passe oublié »
 \item Le système demande l'email de l'utilisateur
 \item Le système envoie un lien de réinitialisation de mot de passe par mail
 \item L'utilisateur suit le lien pour réinitialiser le mot de passe
\end{itemize}\\
\hline
\end{tabular}
\textbf{\caption{Tableau descriptif du cas d'utilisation « S'authentifier »}}
\label{table:tableau descriptif du cas d'utilisation « S'authentifier »}
\end{table}

\vfill

\subsubsection{Diagramme de cas d'utilisation « gérer les commandes d'achats de matière première » :}
Ce cas permet au Responsable Achats de gérer les commandes d'achat de matières premières : consulter, ajouter, modifier et supprimer.\par

\begin{figure}[H]%
 \center%
 \setlength{\fboxsep}{0pt}%
 \setlength{\fboxrule}{0.0pt}%
 \fbox{
 \includegraphics[width=\linewidth]{gerer_achat.png}%
 }
 \caption{Diagramme de cas d'utilisation « gérer les commandes d'achats de matière première »}%
 \end{figure}
 
\textbf{Tableau descriptif :}
\begin{table}[!h]
\renewcommand{\arraystretch}{1.8}
\setlength{\tabcolsep}{16pt}
\centering
\begin{tabular}{|l|p{12cm}|}

\hline
Titre & Gérer les commandes d'achats de matière première \\
\hline
Objectif & Permettre au Responsable Achats de gérer les commandes d'achat de matières premières. \\
\hline
Acteur(s) & Responsable Achats (acteur principal) \\
\hline
Pré-condition & L'utilisateur est authentifié avec le rôle « Responsable Achats ». Au moins un fournisseur actif et une matière première sont référencés dans le système. \\
\hline
Post-condition & La commande d'achat est créée, modifiée ou supprimée selon l'action effectuée. Les stocks sont mis à jour en conséquence. \\
\hline
Scénario nominal & - Le Responsable Achats accède au module « Achats ».\par
 - Le système affiche la liste des commandes d'achat existantes.\par
 - L'utilisateur choisit une action : consulter / Ajouter / Modifier / Supprimer une commande d'achat.\par
 - Pour ajouter : l'utilisateur sélectionne la matière première, le fournisseur, saisit la quantité et la date de livraison.\par
 - Le système vérifie la disponibilité et l'état actif du fournisseur.\par
 - L'utilisateur valide la commande d'achat.\par
 - Le système enregistre la commande en base de données.\par
 - Le système notifie le Responsable Production de la commande.\par
\\
\hline
Scénario alternatif & - Fournisseur inactif ou supprimé : le système affiche un avertissement et propose de choisir un fournisseur alternatif.\par
- Quantité saisie invalide (zéro ou négative) : le système retourne un message d'erreur de validation.\par
\\
\hline
\end{tabular}
\textbf{\caption{Tableau descriptif du cas d'utilisation « Gérer les commandes d'achats de matière première »}}
\label{table:cas d'utilisation « Gérer les commandes d'achats de matière première »}
\end{table}
\vfill
\subsubsection{Diagramme de cas d'utilisation « gérer les produits » : }
Ce cas permet au Magasinier de gérer les articles (CRUD) : matières premières, produits semi-finis et produits finis.\par

\begin{figure}[H]%
 \center%
 \setlength{\fboxsep}{0pt}%
 \setlength{\fboxrule}{0.0pt}%
 \fbox{
 \includegraphics[width=\linewidth]{suivre_stock.png}%
 }
 \caption{Diagramme de cas d'utilisation « gérer les produits »}%
\end{figure}
\textbf{Tableau descriptif :}
\begin{table}[!htbp]
\renewcommand{\arraystretch}{1.8}
\setlength{\tabcolsep}{16pt}
\centering
\begin{tabular}{|l|p{12cm}|}

\hline
Titre & Gérer les produits \\
\hline
Objectif & Permettre au Magasinier de gérer les produits (matières premières, semi-finis, finis).
 \\
\hline
Acteur(s) & Magasinier (acteur principal) \\
\hline
Pré-condition & L'utilisateur est authentifié avec le rôle « Magasinier ». Des données de stock existent dans le système. \\
\hline
Post-condition & Le référentiel des produits est mis à jour. Les mouvements de stock sont enregistrés en base de données.
 \\
\hline
Scénario nominal & \begin{itemize}
 \item Le Magasinier accède au module « Stock ».
\item Le système affiche la liste des produits avec leurs quantités et catégories.
\item L'utilisateur choisit une action : consulter / Ajouter / Modifier / Supprimer un produit.
\item Pour ajouter : l'utilisateur saisit le nom, la catégorie, la quantité et le seuil minimal.
\item Le système valide les données et enregistre le produit.
\item Le système compare les quantités aux seuils minimaux définis.
\item Si un article est en dessous du seuil, le système affiche une alerte visuelle.
\item Le Magasinier peut exporter le rapport de stock si nécessaire.
\end{itemize}
\\
\hline
Scénario alternatif & \begin{itemize}
 \item Stock inférieur au seuil minimal : le système génère une alerte et envoie une notification au Responsable Achats pour déclencher un réapprovisionnement.
\item Aucun article trouvé : le système affiche un message indiquant l'absence de données pour le filtre sélectionné.
\end{itemize} \\
\hline
\end{tabular}
\textbf{\caption{Tableau descriptif du cas d'utilisation « Gérer les produits »}}
\label{table:cas d'utilisation « Gérer les produits »}
\end{table}
\vfill
\subsubsection{Diagramme de cas d'utilisation « gérer les ordres de production » :}
Ce cas permet au Responsable Production de gérer les ordres de production, les produits, les matières premières et les nomenclatures BOM.\par

\begin{figure}[H]%
 \center%
 \setlength{\fboxsep}{0pt}%
 \setlength{\fboxrule}{0.0pt}%
 \fbox{
 \includegraphics[width=\linewidth]{gerer_production.png}%
 }
 \caption{Diagramme de cas d'utilisation « gérer les ordres de production » }%
 \end{figure}
\textbf{Tableau descriptif :}
\begin{table}[!htbp]
\renewcommand{\arraystretch}{1.8}
\setlength{\tabcolsep}{16pt}
\centering
\begin{tabular}{|l|p{12cm}|}

\hline
Titre & Gérer les ordres de production \\
\hline
Objectif & Permettre au Responsable Production de planifier, lancer et suivre les ordres de production et les nomenclatures BOM.
 \\
\hline
Acteur(s) & Responsable Production (acteur principal) \\
\hline
Pré-condition & L'utilisateur est authentifié avec le rôle « Responsable Production ». Au moins un produit avec nomenclature BOM est défini. Les matières premières nécessaires sont référencées en stock. \\
\hline
Post-condition & L'ordre de production est enregistré et/ou démarré. Le stock de matières premières est mis à jour en conséquence. Les parties prenantes sont notifiées.
 \\
\hline
Scénario nominal &
- Le Responsable Production accède au module « Production ».\par
- Le système affiche le tableau de bord de production (planifiées, en cours, terminées).\par
- L'utilisateur crée un nouvel ordre de production (produit, quantité, date souhaitée).\par
- Le système vérifie la disponibilité des matières premières mentionnées dans la BOM.\par
- Si les MP sont disponibles, le système planifie l'ordre et le met en statut « Planifié ».\par
- L'utilisateur lance l'ordre de production.\par
- Le système met à jour le statut en « En cours » et déduit les quantités de MP du stock.\par
- Le système notifie le Magasinier de la consommation de matières premières.\par
- L'utilisateur peut consulter le statut d'avancement à tout moment.\par
- Une fois terminée, l'utilisateur valide la fin de production et le système incrémente le stock de produits finis.\par
\\
\hline
Scénario alternatif & \begin{itemize}
 \item Matières premières insuffisantes : le système affiche une alerte et envoie une notification au Responsable Achats pour un réapprovisionnement urgent. La production ne peut pas démarrer.
 \item Annulation d'un ordre de fabrication planifié : lutilisateur peut annuler un ordre en statut « Planifié ». Les quantités réservées sont restituées au stock.
\end{itemize}\\
\hline
\end{tabular}
\textbf{\caption{Tableau descriptif du cas d'utilisation « Gérer les ordres de production »}}
\label{table:cas d'utilisation « Gérer les ordres de production » }
\end{table}
\vfill
\subsubsection{Diagramme de cas d'utilisation « gérer les commandes clients » :}
Ce cas permet au Responsable Ventes de gérer les commandes clients : le Client passe commande, le Responsable Ventes consulte les produits et valide.\par

\begin{figure}[H]%
 \center%
 \setlength{\fboxsep}{0pt}%
 \setlength{\fboxrule}{0.0pt}%
 \fbox{
 \includegraphics[width=\linewidth]{creer_commande_client.png}%
 }
 \caption{Diagramme du cas d'utilisation « gérer les commandes clients »}%
 \end{figure}

 \textbf{Tableau descriptif :}
 \begin{table}[!htbp]
 \renewcommand{\arraystretch}{1.8}
 \setlength{\tabcolsep}{16pt}
 \centering
 \begin{tabular}{|l|p{12cm}|}
 \hline
 Titre & Gérer les commandes clients \\
 \hline
 Objectif & Permettre au Responsable Ventes de gérer les commandes clients.
 \\
 \hline
 Acteur(s) & Responsable Ventes (acteur principal) \\
 \hline
 Pré-condition & L'utilisateur est authentifié avec le rôle « Responsable Ventes ». Au moins un produit est référencé dans le système. \\
 \hline
 Post-condition & La commande client est enregistrée dans la base de données avec le statut « En cours ». Une notification est envoyée au Responsable Production.
 \\
 \hline
 Scénario nominal & \begin{itemize}
 \item Le Client passe une commande via l'interface.
 \item Le Responsable Ventes accède au module « Ventes ».
 \item Le système affiche la liste des commandes clients.
 \item Le Responsable Ventes consulte les produits disponibles.
 \item Le Responsable Ventes valide la commande client.
 \item Le système enregistre la commande en base de données.
 \item Le système envoie une notification au Responsable Production.
 \item Le système affiche un message de confirmation avec le numéro de commande généré.
 \end{itemize}\\
 \hline
 Scénario alternatif & - Données invalides : le système affiche un message d'erreur précisant les champs incorrects ou manquants. L'utilisateur est invité à corriger les données.\par
 - Produit non disponible : le système propose d'ajouter le produit ou de choisir un article existant.\\
 \hline
 \end{tabular}
 \textbf{\caption{Tableau descriptif du cas d'utilisation « Gérer les commandes clients »}}
 \label{table:cas d'utilisation « Gérer les commandes clients »}
 \end{table}
 \vfill
\subsubsection{Diagramme de cas d'utilisation « gérer les fournisseurs » :}
Ce cas permet au Responsable Achats de gérer les fournisseurs (CRUD). Le système empêche la suppression des fournisseurs liés à des commandes actives.\par

\begin{figure}[H]%
 \center%
 \setlength{\fboxsep}{0pt}%
 \setlength{\fboxrule}{0.0pt}%
 \fbox{
 \includegraphics[width=\linewidth]{gerer_fournisseurs.png}%
 }
 \caption{Diagramme du cas d'utilisation « gérer les fournisseurs »}%
 \end{figure}
 \vspace{3cm}
\textbf{Tableau descriptif :}

 \begin{table}[!htbp]
 \renewcommand{\arraystretch}{1.8}
 \setlength{\tabcolsep}{16pt}
 \centering
 \begin{tabular}{|l|p{12cm}|}

 \hline
 Titre & Gérer les fournisseurs \\
 \hline
 Objectif & Permettre au Responsable Achats de gérer les fournisseurs.
 \\
 \hline
 Acteur(s) & Responsable Achats (acteur principal) \\
 \hline
 Pré-condition & L'utilisateur est authentifié avec le rôle « Responsable Achats ».
 \\
 \hline
 Post-condition & Le référentiel des fournisseurs est mis à jour (ajout, modification ou suppression effectuée).
 \\
 \hline
 Scénario nominal & \begin{itemize}
 \item Le Responsable Achats accède au module Fournisseurs.
 \item Le système affiche la liste des fournisseurs existants.
 \item L'utilisateur choisit une action : ajouter / Modifier / Supprimer / Consulter.
 \item Ajouter: l'utilisateur remplit le formulaire avec les informations du fournisseur (nom, contact, matières livrées, délai moyen).
 \item Modifier: l'utilisateur sélectionne un fournisseur et modifie ses informations.
 \item Supprimer: l'utilisateur sélectionne un fournisseur et confirme la suppression.
 \item Le système valide les données saisies.
 \item Le système enregistre les modifications en base de données.
 \item Le système affiche un message de succès.
 \end{itemize}\\

 \hline
 Scénario alternatif & - Suppression d'un fournisseur lié à des ordres en cours : le système affiche un avertissement et bloque la suppression tant que les ordres sont actifs.\par
 - Informations manquantes (formulaire incomplet) : le système affiche les champs obligatoires manquants et empêche la sauvegarde.\\
 \hline
 \end{tabular}
 \textbf{\caption{Tableau descriptif du cas d'utilisation « Gérer les fournisseurs »}}
 \label{table:cas d'utilisation « Gérer les fournisseurs »}
 \end{table}
 \vfill
\subsubsection{Diagramme de cas d'utilisation « gérer les utilisateurs » :}
Ce cas permet au PDG de gérer les comptes utilisateurs : création, modification, consultation, désactivation et suppression des accès.\par

\begin{figure}[H]%
 \center%
 \setlength{\fboxsep}{0pt}%
 \setlength{\fboxrule}{0.0pt}%
 \fbox{
 \includegraphics[width=\linewidth]{gereutilisateurs.png}%
 }
 \caption{Diagramme du cas d'utilisation « gérer les utilisateurs »}%
 \end{figure}

 \textbf{Tableau descriptif :}

\begin{table}
\renewcommand{\arraystretch}{1.8}
\setlength{\tabcolsep}{16pt}
\centering
\begin{tabular}{|l|p{12cm}|}

\hline
Titre & Gérer les utilisateurs \\
\hline
Objectif & Permettre au PDG de gérer les comptes utilisateurs et d'attribuer les rôles.
 \\
\hline
Acteur(s) & PDG (acteur principal) \\
\hline
Pré-condition & L'utilisateur est authentifié avec le rôle « Administrateur ».
\\
\hline
Post-condition & Le compte utilisateur est créé, modifié, désactivé ou supprimé.\\
\hline
Scénario nominal &\begin{itemize}
 \item Le PDG accède à la rubrique « Gestion des utilisateurs ».
\item Le système affiche la liste des comptes existants avec leurs rôles et statuts.
\item L'utilisateur choisit une action : consulter profil / Créer / Modifier / Désactiver / Supprimer.
\item Le système valide les données (email unique, rôle valide).
\item Le système enregistre les modifications en base de données.
\item Le système affiche une confirmation de l'opération.

\end{itemize}\\
\hline
Scénario alternatif & \begin{itemize}
\item Email déjà utilisé : le système affiche une erreur et empêche la création du doublon.
\item Tentative de suppression d'un compte actif : le système impose la désactivation préalable avant toute suppression.
\item Tentative de suppression du dernier compte administrateur : le système bloque l'opération pour éviter de perdre tout accès d'administration.
\item Rôle incompatible avec des données existantes : si le rôle est modifié alors que l'utilisateur a des actions en cours (ex. ordres de production en attente), le système affiche un avertissement.
\end{itemize}\\
\hline
\end{tabular}
\textbf{\caption{Tableau descriptif du cas d'utilisation « Gérer les utilisateurs »}}
\label{table:cas d'utilisation « Gérer les utilisateurs »}
\end{table}
\vfill
\subsubsection{Diagramme de cas d'utilisation « Gérer les matières premières » :}
Ce cas permet au Magasinier de gérer les matières premières (CRUD) : ajouter, modifier et supprimer.\par

\begin{figure}[H]%
 \center%
 \setlength{\fboxsep}{0pt}%
 \setlength{\fboxrule}{0.0pt}%
 \fbox{
 \includegraphics[width=\linewidth]{gerermp.png}%
 }
 \caption{Diagramme du cas d'utilisation « Gérer les matières premières »}%
 \end{figure}
 \textbf{Tableau descriptif :}
 \begin{table}[h!]
\renewcommand{\arraystretch}{1.8}
\setlength{\tabcolsep}{16pt}
\centering
\begin{tabular}{|l|p{12cm}|}

\hline
Titre & Gérer les matières premières \\
\hline
Objectif & Permettre au Magasinier de gérer les matières premières. \\
\hline
Acteur(s) & Magasinier (acteur principal) \\
\hline
Pré-condition & L'utilisateur est authentifié avec le rôle « Magasinier ». Le référentiel des produits existe.\\
\hline
Post-condition & La matière première est ajoutée, modifiée ou supprimée dans le référentiel.\\
\hline
Scénario nominal & \begin{itemize}
 \item Le Magasinier accède au module « Stock ».
 \item Le système affiche la liste des matières premières existantes.
 \item L'utilisateur choisit une action : ajouter MP / Modifier MP / Supprimer MP.
 \item Pour ajouter : l'utilisateur saisit le nom, l'unité de mesure et le seuil d'alerte.
 \item Le système valide les données et enregistre la matière première.
 \item Le système affiche une confirmation de l'opération.
\end{itemize}
\\
\hline
Scénario alternatif & - Données invalides : le système affiche un message d'erreur et empêche la sauvegarde.\par
- Suppression d'une MP liée à des produits ou commandes actives : le système bloque la suppression et affiche un avertissement.\par
\\
\hline
\end{tabular}
\textbf{\caption{Tableau descriptif du cas d'utilisation « Gérer les matières premières »}}
\label{table:cas d'utilisation « Gérer les matières premières »}
\end{table}

\vfill

\subsubsection{Diagramme de cas d'utilisation « Consulter le Dashboard » :}
Ce cas permet au PDG de consulter un tableau de bord avec les indicateurs des modules ventes, production, achats et stocks, avec filtrage temporel et suggestions IA.\par

\begin{figure}[H]%
 \center%
 \setlength{\fboxsep}{0pt}%
 \setlength{\fboxrule}{0.0pt}%
 \fbox{
 \includegraphics[width=\linewidth]{visualiser_dashboard.png}%
 }
 \caption{Diagramme du cas d'utilisation « consulter le Dashboard »}%
 \end{figure}
 \textbf{Tableau descriptif :}
\begin{table}[!h]
\renewcommand{\arraystretch}{1.8}
\setlength{\tabcolsep}{16pt}
\centering
\begin{tabular}{|l|p{12cm}|}

\hline
Titre & Consulter le Dashboard
 \\
\hline
Objectif & Permettre au PDG de consulter les indicateurs (ventes, production, achats, stocks) et les suggestions IA. \\
\hline
Acteur(s) & PDG (acteur principal) \\
\hline
Pré-condition & L'utilisateur est authentifié avec le rôle « PDG ». Des données existent dans les différents modules (ventes, production, achats, stocks).
\\
\hline
Post-condition & Le PDG a consulté les indicateurs synthétisés et peut prendre des décisions stratégiques basées sur les données affichées.\\
\hline
Scénario nominal & - Le PDG se connecte à \reportAppName{}.\par
- Le système redirige automatiquement vers le tableau de bord (Dashboard).\par
- Le système charge et affiche les indicateurs clés : chiffre d'affaires, commandes en cours, état des stocks, avancement de la production, achats récents.\par
- Le PDG peut filtrer les données par période (semaine, mois, trimestre).\par
- Le PDG navigue entre les différents widgets du tableau de bord.\par
- Le système affiche des alertes visuelles pour les indicateurs critiques (stock bas, retards de production).\par
- Le PDG peut accéder à des suggestions prédictives générées par l'IA.\par
- Le PDG peut accéder en lecture seule aux détails d'un module en cliquant sur un indicateur.\par

\\
\hline
Scénario alternatif & - Aucune donnée disponible pour la période sélectionnée : le système affiche un message indiquant l'absence de données et propose de choisir une autre période.\par
- Indicateur en état critique : le système met en évidence l'indicateur et affiche un résumé de l'anomalie détectée.\par
\\
\hline
\end{tabular}
\textbf{\caption{Tableau descriptif du cas d'utilisation « Consulter le Dashboard »}}
\label{table:cas d'utilisation « Consulter le Dashboard »}
\end{table}

\vfill

\subsection{Diagramme de classe}
\subsubsection{Présentation :}
Le diagramme de classes offre une vue statique de l’application \reportAppName{}. Il définit les
classes, leurs attributs et leurs méthodes, et sert de transition entre les besoins fonctionnels
et le code.
\begin{figure}[H]%
 \center%
 \setlength{\fboxsep}{0pt}%
 \setlength{\fboxrule}{0.0pt}%
 \fbox{
 \includegraphics[width=\linewidth,height=20cm]{class.png}%
 }
 \caption{Diagramme de classe de \reportAppName{}}%
 \end{figure}
 \vfill
\subsection{Diagramme de séquence}
\subsubsection{Présentation :}
Le diagramme de séquence montre les interactions chronologiques entre acteurs, objets et
systèmes dans \reportAppName{}. Il détaille l’ordre des messages échangés et les structures logiques
(boucles, alternatives).
\subsubsection{Diagramme de séquence : « Gérer Fournisseurs » :}
Ce diagramme montre les interactions pour la gestion des fournisseurs (ajout, modification, suppression, consultation).\par
\begin{figure}[H]%
 \center%
 \setlength{\fboxsep}{0pt}%
 \setlength{\fboxrule}{0.0pt}%
 \fbox{
 \includegraphics[width=\linewidth,height=18cm]{seq_gererfournisseurs.png}%
 }
 \caption{Diagramme de séquence : « Gérer Fournisseurs »}%
\end{figure}




\subsubsection{Diagramme de séquence : « Gérer les produits » :
}
Ce diagramme montre la gestion des produits et des nomenclatures BOM par le Responsable Production.\par

\begin{figure}[H]%
 \center%
 \setlength{\fboxsep}{0pt}%
 \setlength{\fboxrule}{0.0pt}%
 \fbox{
 \includegraphics[width=\linewidth,height=18cm]{seq_gererproduitsetbom.png}%
 }
 \caption{Diagramme de séquence : « Gérer les produits »}%
\end{figure}

\subsubsection{Diagramme de séquence : « S'authentifier » :
}
Ce diagramme montre les échanges lors de l'authentification.\par
\begin{figure}[H]%
 \center%
 \setlength{\fboxsep}{0pt}%
 \setlength{\fboxrule}{0.0pt}%
 \fbox{
 \includegraphics[width=\linewidth,height=18cm]{seq_sauthentifier.png}%
 }
 \caption{Diagramme de séquence : « S'authentifier »}%
\end{figure}

\subsubsection{Diagramme de séquence : « Consulter le Dashboard » :
}
Ce diagramme montre la visualisation du Dashboard par le PDG avec les indicateurs et les suggestions IA.\par
\begin{figure}[H]%
 \center%
 \setlength{\fboxsep}{0pt}%
 \setlength{\fboxrule}{0.0pt}%
 \fbox{
 \includegraphics[width=\linewidth,height=18cm]{seq_visualiserdashboard.png}%
 }
 \caption{Diagramme de séquence : « Consulter le Dashboard »}%
\end{figure}

\subsection{Diagramme d'activité}
\subsubsection{Présentation :}
Le diagramme d’activité modélise les flux de travail et les processus de \reportAppName{} en décomposant
chaque activité en étapes logiques. Il répartit les responsabilités par acteur.
\subsubsection{Diagramme d'activité : « Gérer les fournisseurs » :}
Ce diagramme montre la gestion des fournisseurs avec trois parcours : consultation, édition et
suppression. Lors de l’édition, le système vérifie les données. Pour la suppression, il vérifie
l’absence d’ordres actifs.
\begin{figure}[H]%
 \center%
 \setlength{\fboxsep}{0pt}%
 \setlength{\fboxrule}{0.0pt}%
 \fbox{
 \includegraphics[width=\linewidth,height=15cm]{workflow_gererfournisseurs.png}%
 }
 \caption{Diagramme d'activité : « Gérer les fournisseurs »}%
\end{figure}
\subsubsection{Diagramme d'activité : « Gérer les ordres de production » :}
Ce diagramme montre le processus de production : transformation des ordres de fabrication en
produits finis avec mise à jour des stocks. Le système vérifie la disponibilité des composants
via la nomenclature BOM. En cas de pénurie, le processus est interrompu et les Achats sont
alertés. Si les composants sont disponibles, la production est lancée.
\begin{figure}[H]%
 \center%
 \setlength{\fboxsep}{0pt}%
 \setlength{\fboxrule}{0.0pt}%
 \fbox{
 \includegraphics[width=\linewidth,height=15cm]{workflow_gererproduction.png}%
 }
 \caption{Diagramme d'activité : « Gérer les ordres de production »}%
\end{figure}

\subsubsection{Diagramme d'activité : « S'authentifier » :}
Ce diagramme montre le processus d’authentification avec deux niveaux de vérification : champs
non vides et vérification en base. Après plusieurs échecs, le compte est bloqué. En cas de
succès, un jeton JWT est généré et l’utilisateur est redirigé vers son tableau de bord.
\begin{figure}[H]%
 \center%
 \setlength{\fboxsep}{0pt}%
 \setlength{\fboxrule}{0.0pt}%
 \fbox{
 \includegraphics[width=\linewidth,height=15cm]{workflow_sauthentifier.png}%
 }
 \caption{Diagramme d'activité : « S'authentifier »}%
\end{figure}


{\Large\bfseries Conclusion}\\
Dans ce chapitre, nous avons présenté la conception de \reportAppName{} en utilisant la langage UML. Les diagrammes de cas d’utilisation
ont modélisé les interactions, le diagramme de classes a défini la structure des données, et
les diagrammes de séquence et d’activité ont détaillé les processus. Ces modèles nous guideront
par la suite  à la phase d’implémentation.